\section{Introduction}

Emerging economies now account for a large share of the world's patenting: China's patent office alone received
49\% of the 3.7 million patent applications filed worldwide in 2024, up from 35\% in 2014 \citep{WIPO2025}. Yet
their patents are widely thought to be worth less. Renewal decisions offer a direct test. Owners pay a fee every year to keep a patent in force and stop paying
once it is no longer worth the fee, so the age at which patents lapse reveals what they are worth to their owners
\citep{PakesSchankerman1984,Pakes1986,SchankermanPakes1986}. By this measure the gap looks large. In the Chinese
patent office, patents filed in 2000--2005 by domestic firms were 13 percentage points less likely than those of
foreign owners to be alive eight years after grant.

Two very different stories fit such a gap. In the first, the \emph{invention} is the problem: emerging-market
patents protect weaker technology, and owners rationally drop them sooner. In the second, the \emph{inventor} is
the problem: owners lack the finance, production capacity or market access to turn valuable inventions into
returns, so good patents lapse for want of an owner able to use them. The policy response differ. The first calls
for investment in research capability and links to frontier technology; the second for credit, fee relief or
markets in which patents can find better owners.

This paper separates the two using the legal-event records in PATSTAT, which track renewals and lapses at the
offices that report them. I first establish which offices do: of the large emerging economies, only China, Korea,
Russia, Taiwan and the Czech Republic record lapses consistently and have enough domestic firm
patents.\footnote{India, Mexico, South Africa and most of Latin America, Africa and South-East Asia report no lapses
to PATSTAT; Brazil and Hungary report them but have too few domestic firm patents (Appendix~\ref{app:data}).}
I compare these against five advanced economies (the United States, Japan, Germany, the United Kingdom, and France) using a dataset of approximately three million granted patents held by firms and universities, complemented with the historical renewal-fee schedules compiled from official patent-office sources. The design compares domestic and foreign owners in the same office (exposed to the same fees, law and
enforcement), and asks whether the domestic shortfall is larger in emerging-market offices. The paper makes four
points.

First, the raw shortfall reflects what is patented. Domestic patents in emerging-market offices are cited far less
than foreign ones, and once citations, technology field and owner type are held fixed, the extra domestic shortfall in
emerging-market offices disappears: larger emerging-market firms keep their patents longer than comparable firms in
the advanced economies, and small firms about as long, within a few points either way depending on the quality
controls and on which advanced economies form the benchmark (Section~\ref{sec:survival}). The result holds in China and Korea, but imprecise in the Czech Republic, and
fails for small firms in Russia and, more weakly, in Taiwan.

Second, there is no sign that emerging-market owners are held back by finance or by the cost of renewal
(Section~\ref{sec:constraints}). Small firms everywhere keep the home patent of an invention longer than its
foreign patents, and more so in the emerging economies, the opposite of what owners unable to use their home patents
would do. Credit crunches do not raise small domestic firms' lapses in the emerging economies. And small firms respond only weakly to large fee cuts.

Third, structural renewal models turn these decisions into money (Section~\ref{sec:values}). Valuing each patent from its own renewal history,
without reference to its owner, the home premium (the value of domestic relative to foreign owners' patents in the
same office) is as large in the emerging-market offices as in the advanced ones for all owners together and for
mid-sized and large firms. For small firms it is smaller against the advanced-economy pool, which Japan dominates;
office by office, small firms' premium in China and Korea is as large as in the United States and Germany. The results
hold when every eligible owner is assumed to pay discounted fees provided by Chinese and Korean patent office.

Fourth, two differences remain. Patents sold by small emerging-market firms survive much longer afterwards than
patents sold by small firms in the advanced economies, which suggests that owner capacity matters at the margin and that markets for
technology move patents to better owners. And university patents in China are the one robust weak spot: they lapse
sooner and are worth well under half as much as foreign universities' patents at equal quality. Here the owner
matters too: once sold, Chinese university patents survive 45 points longer, while a sale does not raise the survival of
university patents in the advanced economies. This implies that Chinese universities hold patents that firms find worth keeping but
that they do not use themselves.

I close by pricing the policy levers on a common scale (Section~\ref{sec:stake}). Had every eligible owner taken up
China's and Korea's fee discounts, the returns from a cohort of small firms' patents would barely have changed. Had patents held by small firms generated the same returns as those realized by their buyers, aggregate returns would have increased by up to \$250 million per cohort in China and \$80 million in Korea --- this rises to \$1.1 billion and \$143 million for all domestic owners, universities included, in China and Korea respectively. Similarly, if Chinese universities' patents had generated returns comparable to those of foreign universities operating in China, the gains would have ranged from \$170 million to \$310 million. These are private returns, not welfare, but the ordering is clear: the gains lie in markets for
technology and in what universities patent, not in cheaper renewal.

The paper contributes to three literatures. The first estimates the value of patent rights from renewal decisions
\citep{PakesSchankerman1984,Pakes1986,SchankermanPakes1986,Lanjouw1998,Bessen2008}; I extend it to emerging-market
offices, with fee histories for ten offices over 25 years. The second studies indicators of patent value
\citep{HarhoffSchererVopel2003,HallJaffeTrajtenberg2005,KoganEtAl2017}; I value each patent from its own renewal
path rather than from its owner's characteristics. The third asks whether Chinese and other emerging-market patenting
is of low quality \citep{HuJefferson2009,LiegsalzWagner2013,DangMotohashi2015}; the answer here is that, for firms,
lower survival reflects the technology, not the owners' ability to use their patents. The sale premium connects to
work on markets for patents \citep{Serrano2010,GalassoSchankermanSerrano2013}.

Section~\ref{sec:data} describes the data and design. Section~\ref{sec:survival} compares survival at equal
quality, Section~\ref{sec:constraints} tests for owner constraints, Section~\ref{sec:values} turns renewal decisions
into money, and Section~\ref{sec:stake} compares policy levers. Section~\ref{sec:conclusion} concludes.
